% SPDX-License-Identifier: LPPL-1.3c
% Copyright (C) 2026 Christophe Jorssen
% Author and Current Maintainer: Christophe Jorssen <christophe.jorssen@gmail.com>
% LPPL maintenance status: maintained
% This file is part of luafemm. See LICENSE and LICENSES.md for its terms.
\section{Charged bodies and ideal capacitors}
\label{sec:electric-classics}

The common electric views and profiles also accept exact symmetry models.
Each picture contains one model in a homogeneous, linear, isotropic medium
with $\epsilon=\epsilon_0\epsilon_r$. Spherical models describe a central plane
through a three-dimensional sphere; cylindrical models describe a transverse
plane of an infinite cylinder; the plate model describes two infinite parallel
planes. Their electric field lies in the displayed plane. These are analytical
solutions, with no finite-element mesh, finite-length correction or edge effect.

\subsection{Physical declarations}
\label{sec:electric-classic-model}
The following keys belong to |/femm/electrostatic| and are set before the
picture begins, usually through |femm/electrostatic={...}|. They retain normal
TeX grouping and can be assembled into ordinary PGF styles. Numeric inputs
accept finite decimal literals, including scientific notation. Lengths and
centres use the model's |length unit|; densities, voltages and returned field
values use SI. Source density and capacitor voltage are required physical
inputs; zero and negative values are allowed.

\begin{key}{/femm/electrostatic/model=\meta{choice} (initially point charges)}
Choices are:
\begin{itemize}
\item |point charges| (the default);
\item |charged sphere| or |charged cylinder|;
\item |parallel plate capacitor|, |cylindrical capacitor| or |spherical capacitor|;
\item |finite elements| (the planar/axisymmetric Poisson model of
      Section~\ref{sec:electric-fe});
\item |source correction| (singular line/axial point sources with a regular
      correction, Section~\ref{sec:electric-hybrid}).
\end{itemize}
The first accepts the point-charge nodes of Section~\ref{sec:electric-charge}.
Each classical choice creates one complete analytical configuration and rejects
additional physical point-charge nodes. Ordinary TikZ annotations are allowed.
Combining multiple ideal capacitors or adding arbitrary conductors is not
supported. No magnetic region, boundary or circuit declaration applies here.
\end{key}

\begin{key}{/femm/electrostatic/center x=\meta{number} (initially 0)}
Physical centre's $x$ coordinate in model units. For plates this is the centre
of the visible pair and the origin of the signed normal coordinate.
\end{key}
\begin{key}{/femm/electrostatic/center y=\meta{number} (initially 0)}
Physical centre's $y$ coordinate in model units. Centres belong to the root
physical frame; a later local scope or pic shift does not reposition the model.
The root picture transform changes the view of all its data consistently.
\end{key}
\begin{key}{/femm/electrostatic/radius=\meta{positive number} (initially 1)}
Radius $R$ of a charged sphere or cylinder, in model units. It is irrelevant
to a capacitor, whose electrode radii or separation are specified separately.
\end{key}
\begin{key}{/femm/electrostatic/charge distribution=\meta{choice} (initially volume)}
For a charged sphere or cylinder, |volume| prescribes a uniform volume density
and |surface| a uniform density on the spherical or cylindrical surface.
The interior and exterior have the same permittivity. A volume charge is a
prescribed distribution, not a conductor in electrostatic equilibrium.
\end{key}
\begin{key}{/femm/electrostatic/volume charge density=\meta{number} (initially \normalfont empty)}
Signed $\rho$ in $\mathrm{C/m^3}$, required for a |volume| distribution.
Omitting it fails. Providing both density keys, or a density on a capacitor,
is an error. It is not scaled when |length unit| changes.
For |model=finite elements|, this key instead sets the background volume density;
empty means zero. Physical region paths may replace it locally.
\end{key}
\begin{key}{/femm/electrostatic/surface charge density=\meta{number} (initially \normalfont empty)}
Signed $\sigma$ in $\mathrm{C/m^2}$, required for a |surface| distribution.
The normal displacement jump is $D_{r,\mathrm{out}}-D_{r,\mathrm{in}}=\sigma$.
Only the selected density key may be supplied.
\end{key}
\begin{key}{/femm/electrostatic/reference radius=\meta{positive number} (initially \normalfont empty)}
Required for a charged infinite cylinder: $V(r_0)$ equals |reference potential|
at this positive radius $r_0$, in model units. It may lie inside or outside the
charged cylinder. The exterior logarithmic potential has no finite reference
at infinity. For planar |source correction|, empty instead selects one model
unit as the arbitrary logarithmic source reference; changing it preserves
the total prescribed-voltage solution. It is rejected for axial source
correction, spheres, capacitors and ordinary point charges.
\end{key}
\begin{key}{/femm/electrostatic/reference potential=\meta{number} (initially 0)}
Potential in volts at infinity for a charged sphere, at |reference radius| for
a charged cylinder, on the outer electrode of a radial capacitor, or on the
positive-normal electrode of the plate capacitor. It shifts $V$ without
changing $\mathbf E$, $\mathbf D$, charge or energy. The point-charge model
keeps its fixed zero reference at infinity; this key does not configure it.
\end{key}
\begin{key}{/femm/electrostatic/voltage=\meta{number} (initially \normalfont empty)}
Required signed electrode voltage difference $U$ in volts for a capacitor.
For radial capacitors $V(a)=V_0+U$, $V(b)=V_0$; for plates
$V(\xi=-d/2)=V_0+U$, $V(\xi=d/2)=V_0$. A positive $U$ produces an outward
radial field or a positive-normal field. Zero gives a constant potential and
no electric field; negative values reverse the field. This key is rejected on
prescribed charge models.
\end{key}
\begin{key}{/femm/electrostatic/inner radius=\meta{positive number} (initially .7)}
Inner electrode radius $a$ of a cylindrical or spherical capacitor, in model
units. It must be strictly smaller than |outer radius|. The inner region is
an equipotential conductor with zero field.
\end{key}
\begin{key}{/femm/electrostatic/outer radius=\meta{positive number} (initially 1.4)}
Outer electrode radius $b$ of a cylindrical or spherical capacitor, in model
units. Outside it $V=V_0$ and $\mathbf E=\mathbf 0$. Electrode thickness is not
modelled; the outer electrode is represented by its boundary surface.
\end{key}
\begin{key}{/femm/electrostatic/separation=\meta{positive number} (initially 1)}
Plate separation $d$ in model units. The two planes lie at signed normal
coordinates $-d/2$ and $d/2$ relative to the centre. The field is zero outside
the interval, including both ideal conductor sides.
\end{key}
\begin{key}{/femm/electrostatic/normal angle=\meta{number} (initially 0)}
Plate normal angle $\theta$ in degrees, measured counterclockwise from physical
$x$. With $\mathbf n=(\cos\theta,\sin\theta)$ and centre $\mathbf c$,
$\xi=(\mathbf r-\mathbf c)\cdot\mathbf n$. This changes the physical orientation
of the plates; the picture's TikZ transform only changes their presentation.
\end{key}
\begin{key}{/femm/electrostatic/interface value=\meta{choice} (initially upper)}
Choose |lower| or |upper| for a sample exactly on a charged surface or electrode.
These denote limits from smaller or larger radial coordinate, or signed normal
coordinate for plates. They do not denote lower/higher potential or the sign
of the field. A relative geometric equality tolerance is used, bounded by the
electrode spacing for a thin radial gap. $V$ is continuous and unaffected.
Profiles that cross an interface explicitly retain both limits, independently
of this default; profiles running along an interface use this choice.
\end{key}

\begin{codeexample}[width=0pt]
\begin{tikzpicture}[femm/electrostatic={model=parallel plate capacitor,
    center x=.3,center y=.2,separation=.8,normal angle=30,
    voltage=-60,reference potential=10,interface value=lower,
    xmin=-1.3,xmax=1.9,ymin=-1.2,ymax=1.6}]
  \pic[femm/electric={plate span=4}] {femm electric geometry};
  \pic[femm/electric={lines per charge=7}] {femm electric field};
  \pic[femm/electric={levels={-40,-20,0}}] {femm equipotentials};
  \node[above] at (.3,.2) {$U=-60\,\mathrm V$};
\end{tikzpicture}
\end{codeexample}

A named |snapshot| owns all these physical parameters. Restore it without
repeating or overriding |model|, radii, densities, voltage, reference or
interface convention: explicit assignments are rejected even if equal to the
saved value. Its original |length unit| and |relative permittivity| must still
be repeated when nondefault. View bounds, transforms and rendering keys can
change. Profiles own their captured model and remain valid after the picture
ends. No persistent file or FEMM-native interchange is provided for these
analytical configurations.

\subsection{Geometry, field lines and discontinuities}
\label{sec:electric-classic-views}
\begin{key}{/tikz/pics/femm electric geometry=\meta{pic}}
Draw the selected classical model from its physical parameters. A charged body
uses its physical radius; a radial capacitor draws both electrode boundaries
and the inner conductor; a plate capacitor draws two straight segments. The
pic adds no physical sources. Place it before the field if its fill should be
behind the lines. It is an error to use this pic with |point charges|.
Circle boundaries are drawn as 180 straight pieces; field values and profile
intersections use the exact circles independently of this display approximation.
\end{key}
\begin{key}{/femm/electric/plate span=\meta{positive number} (initially 4)}
Length of each visible plate segment, in model units, for the geometry pic.
This is a display extent only: it never makes the physical plate finite.
It is independent of the field window and the separation. Choose a span wider
than the window's tangent extent when showing the infinite-plate approximation.
\end{key}
\begin{key}{/tikz/femm charged body=\meta{style}}
Default style for a prescribed uniform volume charge: red outline, pale red
fill, .6\,pt line width. Use ordinary |.append style| to override drawing keys.
\end{key}
\begin{key}{/tikz/femm charged surface=\meta{style}}
Default style for a prescribed uniform surface charge: red outline, .9\,pt
line width, no fill. It does not create a magnetic region or a conductor law.
\end{key}
\begin{key}{/tikz/femm electric conductor=\meta{style}}
Default inner-conductor style: dark grey outline, pale grey fill, .8\,pt width.
Changing the fill does not change the physical permittivity or electrode voltage.
\end{key}
\begin{key}{/tikz/femm electric electrode=\meta{style}}
Default boundary style for plates and the outer radial electrode: dark grey
outline, .8\,pt width. There is no automatically assigned physical thickness.
\end{key}

The ordinary |femm electric field|, |femm equipotentials| and
|femm electric vectors| pics, and the |femm/electric| family, work with these
models. Classical field lines are exact straight radial or normal segments.
They stop at the view window, a zero-field region, the |null field| threshold,
or |max length|, as appropriate. Their arrows always follow $\mathbf E$.
No integration is necessary: |step|, |tolerance| and |max steps| affect
point-charge integration but do not refine these straight lines.

Without explicit seeds, |lines per charge| is the total number of radial rays
or normal lines for the single classical model. |seed sign| filters by density
or voltage sign; |auto| accepts either. Explicit |seeds| and captured seed
contours select rays/normal lines through those points. Seeds outside the
view or in a zero-field region are skipped. |direction| selects either or both
sides of a seed, while |max length| bounds each traced side in model units.
Uniform angular spacing in a section does not encode three-dimensional flux.

Equipotentials use the common sampled contour renderer. Refine
|contour resolution| for curvature and narrow gaps; it cannot change the
analytical field. A constant-potential region is not a unique isoline.
Avoid requesting its value if a specific contour curve is intended. Vector
arrows that would cross an interface with a field jump are omitted rather than
being drawn through it. Their magnitude convention remains |vector mode|.

Profiles insert exact intersections of the flattened measurement segments with
circles or plates, even when |samples| would miss them. A volume-charge boundary
adds a continuous sample. A genuine field jump adds two one-sided values and
an intervening NaN for E/D components and norms, while V remains finite. Thus
PGFPlots leaves the field jump open and keeps the potential continuous. This
also works on reversed contours and through both sides of a radial model.
The number of rows can exceed |samples|. Curved measurement paths still have
the documented geometric flattening tolerance; a very narrow missed curve
excursion requires refining that tolerance.

\subsection{Integral quantities and formatting}
\label{sec:electric-classic-values}
\begin{command}{\femmelectricvalue\oarg{options}}
Print an electric integral quantity using PGF's number formatter.
The FE electrode selector and additional quantities are documented in
Section~\ref{sec:electric-fe-views}. Options
belong to |/femm/electric value| and are local to this invocation. It works
inside a picture or after restoring a named model, in all three formats;
PGFPlots is not required. The result is formatted text, not an expandable
numeric macro. Supply the unit in surrounding TeX. Unknown models, unavailable
quantities and nonfinite numerical results are errors.
\end{command}
\begin{key}{/femm/electric value/model=\meta{name} (initially \normalfont empty)}
Use a saved electric snapshot. Empty means the current electric picture;
outside one, a name is required. Snapshot names are document-wide.
\end{key}
\begin{key}{/femm/electric value/quantity=\meta{choice} (initially charge)}
Analytical choices are:
\begin{itemize}
\item |charge|, |charge per length| or |charge per area|;
\item |capacitance|, |capacitance per length| or |capacitance per area|;
\item |energy|, |energy per length| or |energy per area|;
\item |voltage|.
\end{itemize}
Only physically appropriate normalizations are available, as listed below.
An infinite cylinder does not have a finite total charge or capacitance.
Additional FE diagnostics are listed in Section~\ref{sec:electric-fe-views}.
\end{key}
\begin{key}{/femm/electric value/number format=\marg{PGF options} (initially \{sci,precision=3\})}
Native |/pgf/number format| settings. For example |fixed,precision=1| prints
one fixed decimal place when rounding needs it. Formatting does not modify
the underlying SI quantity. Use PGF's usual precision and zero-fill controls.
\end{key}

\begin{center}\small
\begin{tabular}{lll}
Model & Available quantities & Units\\\hline
Point charges; charged sphere & charge & C\\
Charged infinite cylinder & charge per length & C/m\\
Spherical capacitor & capacitance, charge, energy & F, C, J\\
Cylindrical capacitor & the same, each |per length| & F/m, C/m, J/m\\
Parallel plate capacitor & the same, each |per area| & F/m$^2$, C/m$^2$, J/m$^2$\\
Any capacitor & voltage & V
\end{tabular}
\end{center}
The capacitor charge is the signed charge of the inner electrode, or of the
negative-normal plate. The opposite electrode has the opposite charge. Energy
is the nonnegative stored energy. No self-energy of a point-charge or prescribed
charged-body model is returned.

\begin{codeexample}[width=0pt]
\begin{tikzpicture}[femm/electrostatic={model=spherical capacitor,
    length unit=mm,relative permittivity=2,
    inner radius=7,outer radius=14,voltage=80,
    xmin=-16,xmax=16,ymin=-16,ymax=16},x=1mm,y=1mm,
    femm electric conductor/.append style={fill=yellow!15},
    femm electric electrode/.append style={draw=blue!70!black,thick}]
  \pic {femm electric geometry};
  \pic[femm/electric={lines per charge=8}] {femm electric field};
  \node[below] at (0,-16) {$U=\femmelectricvalue[quantity=voltage,
    number format={fixed,precision=0}]\,\mathrm V$};
  \node[above] at (0,16)
    {$C=\femmelectricvalue[quantity=capacitance]\,\mathrm F$};
\end{tikzpicture}
\end{codeexample}

\subsection{Equations and validation limits}
\label{sec:electric-classic-methods}
Here all lengths in the equations are in metres. For a sphere of radius $R$,
$Q=4\pi R^3\rho/3$ or $Q=4\pi R^2\sigma$. Outside it,
$E_r=Q/(4\pi\epsilon r^2)$ and $V=V_0+Q/(4\pi\epsilon r)$.
Inside a volume charge,
\[
 E_r=\frac{\rho r}{3\epsilon},\qquad
 V=V_0+\frac{\rho(3R^2-r^2)}{6\epsilon}.
\]
Inside a surface charge, $E_r=0$ and $V=V_0+\sigma R/\epsilon$.

For an infinite cylinder, $\lambda=\pi R^2\rho$ or $2\pi R\sigma$.
Define $A=\lambda/(2\pi\epsilon)$ and a continuous potential function
\[
 f(r)=\begin{cases}
 A(1-r^2/R^2)/2,&r<R\text{, volume charge},\\
 0,&r<R\text{, surface charge},\\
 -A\log(r/R),&r\ge R.
 \end{cases}
\]
Then $V=V_0+f(r)-f(r_0)$. Inside a volume charge $E_r=\rho r/(2\epsilon)$;
inside a shell it is zero; outside either distribution $E_r=A/r$.
Thus a reference change cancels exactly from the field and from voltage differences.

Inside the capacitor gap the solutions and capacitances are
\[
\begin{array}{cccc}
 & V-V_0 & \text{field along the normal} & \text{capacitance}\\\hline
\text{plates} & U(1/2-\xi/d) & U/d & C/A=\epsilon/d\\
\text{cylinders} & U\log(b/r)/\log(b/a) & U/[r\log(b/a)]
 & C/L=2\pi\epsilon/\log(b/a)\\
\text{spheres} & Uab(1/r-1/b)/(b-a) & Uab/[(b-a)r^2]
 & C=4\pi\epsilon ab/(b-a)
\end{array}
\]
At either side the potential continues with its electrode value and the field
is zero. In each normalization $Q=CU$ and $W=CU^2/2$. A stable logarithmic ratio
and factored spherical expressions avoid subtracting nearly equal reciprocals
in thin gaps; geometries closer than floating-point resolution remain invalid.

The numerical checks exercise Gauss flux, $\mathbf E=-\nabla V$, continuity of
$V$, displacement jumps, signed voltages, zero sources, units, independent
energy integration, thin gaps and interface-aware profiles. They establish
these symmetry formulas and their rendering contracts, not arbitrary conductor
solutions. There are no dielectric interfaces, finite ends, fringe fields,
floating electrodes, external electric circuits or electrostatic FEMM files
in this analytical interface.
